\chapter{Package and Version Reference}
\label{app:packages}

Agent Framework publishes roughly three dozen NuGet packages under the
\code{MicrosoftAgentFramework} account. This appendix lists them, grouped the way
they actually ship.

\begin{versionbox}
Every version in this appendix was read from \code{nuget.org} in \mafdate{}.
Treat it as a dated snapshot. Before starting new work, check
\code{nuget.org/profiles/MicrosoftAgentFramework} and
\code{github.com/microsoft/agent-framework/releases}.
\end{versionbox}

% ============================================================
\section{Release trains}
\label{sec:trains}
\index{release trains}

The single most useful fact about the package family is that it does not move on
one version line. It moves in \emph{trains} --- groups of packages published
together on the same date with the same version. At the time of writing four were
observable:

\begin{center}
\small
\begin{tabularx}{\linewidth}{@{}llX@{}}
\toprule
\textbf{Train} & \textbf{Published} & \textbf{Contents} \\
\midrule
\mafcore{} & 10 June 2026 & Core, providers, workflows, DevUI, A2A, AG-UI, Foundry client, GitHub Copilot \\
\mafext{} & 22 May 2026 & Harness, Hyperlight/CodeAct, DurableTask, Copilot Studio, Cosmos NoSQL, Foundry hosting, Purview, Azure Functions hosting, Aspire DevUI \\
\code{1.0.0-rc5} & 31 March 2026 & \pkg{Microsoft.Agents.AI.AzureAI} and its declarative companion --- effectively superseded by the \code{Foundry} packages \\
\code{1.0.0-preview} & 31 March 2026 & \pkg{Microsoft.Agents.AI.FoundryMemory} \\
\bottomrule
\end{tabularx}
\end{center}

Three consequences.

\textbf{A single MSBuild property will not work.} The tempting
\code{<PackageVersion Include="..." Version="\$(MafVersion)" />} pattern breaks
the moment you take a dependency from a second train. Pin each package.

\textbf{Version number does not tell you maturity.} \pkg{Microsoft.Agents.AI} at
\mafcore{} is stable; \pkg{Microsoft.Agents.AI.Hosting} at
\code{\mafcore{}-preview.260610.1} is not. Same train, same number, different
suffix. You must read the suffix, not the version.

\textbf{A stale train is a signal.} Packages still on a March \code{rc5} while
the family has shipped twice since are being deprecated in place. The
\code{AzureAI} packages are the example: the \code{Foundry}-named packages are
their continuation.

% ============================================================
\section{Core --- stable}

These carry no prerelease suffix and are covered by the GA support commitment.
Everything in Parts~I and~II of this book uses only these.

\begin{center}
\small
\begin{tabularx}{\linewidth}{@{}llX@{}}
\toprule
\textbf{Package} & \textbf{Version} & \textbf{Purpose} \\
\midrule
\pkg{Microsoft.Agents.AI.Abstractions} & \mafcore{} & Interfaces and abstractions. Reference directly when writing a library that should not pull in the implementation. \\
\pkg{Microsoft.Agents.AI} & \mafcore{} & Core functionality. \api{AIAgent}, \api{ChatClientAgent}, threads, tools, context providers, middleware. \\
\pkg{Microsoft.Agents.AI.OpenAI} & \mafcore{} & OpenAI and Azure OpenAI support. \\
\pkg{Microsoft.Agents.AI.Workflows} & \mafcore{} & Executors, edges, workflow builder, checkpointing. \\
\pkg{Microsoft.Agents.AI.Workflows.Generators} & \mafcore{} & Roslyn source generators giving compile-time route configuration for executors. \\
\pkg{Microsoft.Agents.AI.Workflows.Declarative} & \mafcore{} & YAML-defined workflows. \\
\bottomrule
\end{tabularx}
\end{center}

\begin{note}
Note what is \emph{not} here: any Ollama, Bedrock or Gemini package. Those
providers are reached through the \pkg{Microsoft.Extensions.AI.*} connector
family instead (\S\ref{sec:supporting}), because the provider abstraction lives
below Agent Framework. See \S\ref{sec:seam} in Chapter~\ref{ch:first-agent}.
\end{note}

% ============================================================
\section{Providers and protocols --- prerelease}

\begin{center}
\small
\begin{tabularx}{\linewidth}{@{}llX@{}}
\toprule
\textbf{Package} & \textbf{Version} & \textbf{Purpose} \\
\midrule
\pkg{Microsoft.Agents.AI.Anthropic} & \mafcore{}-preview & Anthropic agents. \\
\pkg{Microsoft.Agents.AI.GitHub.Copilot} & \mafcore{}-rc1 & GitHub Copilot SDK as an agent backend. \\
\pkg{Microsoft.Agents.AI.Foundry} & \mafcore{}-preview & Foundry Agents client. \\
\pkg{Microsoft.Agents.AI.AzureAI} & 1.0.0-rc5 & Foundry Agents, older naming. Superseded. \\
\pkg{Microsoft.Agents.AI.AzureAI.Persistent} & \mafcore{}-preview & Azure AI Persistent Agents. \\
\pkg{Microsoft.Agents.AI.CopilotStudio} & \mafext{}-preview & Copilot Studio integration. \\
\pkg{Microsoft.Agents.AI.A2A} & \mafcore{}-preview & Agent-to-Agent protocol client. \\
\pkg{Microsoft.Agents.AI.AGUI} & \mafcore{}-preview & Agent-User Interaction protocol client. \\
\bottomrule
\end{tabularx}
\end{center}

% ============================================================
\section{Hosting --- prerelease}

\begin{center}
\small
\begin{tabularx}{\linewidth}{@{}llX@{}}
\toprule
\textbf{Package} & \textbf{Version} & \textbf{Purpose} \\
\midrule
\pkg{Microsoft.Agents.AI.Hosting} & \mafcore{}-preview & General agent hosting. \\
\pkg{Microsoft.Agents.AI.Hosting.OpenAI} & \mafcore{}-alpha & OpenAI Responses protocol hosting. Alpha --- expect breaking changes. \\
\pkg{Microsoft.Agents.AI.Hosting.A2A} & \mafcore{}-preview & Hosting A2A agents. \\
\pkg{Microsoft.Agents.AI.Hosting.A2A.AspNetCore} & \mafcore{}-preview & A2A hosting in ASP.NET Core. \\
\pkg{Microsoft.Agents.AI.Hosting.AGUI.AspNetCore} & \mafcore{}-preview & AG-UI hosting in ASP.NET Core. \\
\pkg{Microsoft.Agents.AI.Hosting.AzureFunctions} & \mafext{}-preview & Durable agent hosting on Azure Functions. \\
\pkg{Microsoft.Agents.AI.Foundry.Hosting} & \mafext{}-preview & Foundry Hosted Agents. Chapter~\ref{ch:hosting}. \\
\bottomrule
\end{tabularx}
\end{center}

\begin{warning}
\index{namespaces!versus package names}
The namespace and the package name are not always the same word.
\code{using Microsoft.Agents.AI.Foundry.Hosting;} comes from the package of that
name --- but it sits on the \mafext{} train, not the \mafcore{} one, so adding it
to a solution otherwise pinned at \mafcore{} means adding a second version number.
This is the most common cause of a confusing restore failure.
\end{warning}

% ============================================================
\section{Execution, storage and governance --- prerelease}

\begin{center}
\small
\begin{tabularx}{\linewidth}{@{}llX@{}}
\toprule
\textbf{Package} & \textbf{Version} & \textbf{Purpose} \\
\midrule
\pkg{Microsoft.Agents.AI.Harness} & \mafext{}-preview & \api{HarnessAgent}: a pre-configured agent for long-running tasks. Chapter~\ref{ch:harness}. \\
\pkg{Microsoft.Agents.AI.Hyperlight} & \mafext{}-preview & Hyperlight-backed CodeAct --- sandboxed execution of model-generated code. \\
\pkg{Microsoft.Agents.AI.DurableTask} & \mafext{}-preview & Distributed durable execution. Chapter~\ref{ch:hosting}. \\
\pkg{Microsoft.Agents.AI.CosmosNoSql} & \mafext{}-preview & Cosmos DB NoSQL implementations of \api{ChatHistoryProvider} and \api{CheckpointStore}. \\
\pkg{Microsoft.Agents.AI.FoundryMemory} & 1.0.0-preview & Foundry Memory integration. Chapter~\ref{ch:memory}. \\
\pkg{Microsoft.Agents.AI.Purview} & \mafext{}-rc1 & Connecting generative AI apps to Microsoft Purview. \\
\pkg{Microsoft.Agents.AI.Declarative} & \mafcore{}-rc1 & Declarative \emph{agents}, as distinct from declarative workflows. \\
\bottomrule
\end{tabularx}
\end{center}

\begin{note}
\pkg{Microsoft.Agents.AI.CosmosNoSql} is worth noting even if you never use
Cosmos DB, because its description names the two storage abstractions the
framework expects you to implement: \api{ChatHistoryProvider} and
\api{CheckpointStore}. Chapter~\ref{ch:memory} implements the first against
PostgreSQL; Chapter~\ref{ch:workflows-advanced} uses the second.
\end{note}

% ============================================================
\section{Tooling}

\begin{center}
\small
\begin{tabularx}{\linewidth}{@{}llX@{}}
\toprule
\textbf{Package} & \textbf{Version} & \textbf{Purpose} \\
\midrule
\pkg{Microsoft.Agents.AI.DevUI} & \mafcore{}-preview & Developer UI for inspecting agents. Chapter~\ref{ch:observability}. \\
\pkg{Aspire.Hosting.AgentFramework.DevUI} & \mafext{}-preview & DevUI as an Aspire resource. \\
\pkg{Microsoft.Agents.AI.ProjectTemplates} & 1.3.0-preview & \code{dotnet new} templates: an AI Agent Web API project. \\
\pkg{Microsoft.Extensions.AI.Templates} & --- & Sibling package: AI chat web app templates with vector storage and ingestion. Not published by this account. \\
\pkg{Microsoft.McpServer.ProjectTemplates} & --- & Sibling package: an MCP server template. Relevant to Chapter~\ref{ch:tools}. \\
\bottomrule
\end{tabularx}
\end{center}

Installing the templates is worth doing before Chapter~\ref{ch:first-agent}, if
only to see what Microsoft considers a default project shape.
\S\ref{sec:templates} covers what they produce and why the book scaffolds by hand
regardless:

\begin{shellcmd}
dotnet new install Microsoft.Agents.AI.ProjectTemplates::1.3.0-preview.1.26251.3
dotnet new list agent
\end{shellcmd}

% ============================================================
\section{Supporting packages}
\label{sec:supporting}

Not published by the Agent Framework account, but required by the samples.

\begin{center}
\small
\begin{tabularx}{\linewidth}{@{}lX@{}}
\toprule
\textbf{Package} & \textbf{Purpose} \\
\midrule
\pkg{Microsoft.Extensions.AI} & \api{IChatClient} and the provider abstraction. The layer below Agent Framework. \\
\pkg{Microsoft.Extensions.AI.OpenAI} & \api{AsIChatClient} for OpenAI-family clients. \\
\pkg{Microsoft.Extensions.AI.Ollama} & \api{OllamaChatClient}. Prerelease. Used throughout Parts~I and~II. \\
\pkg{OllamaSharp} & Community alternative providing \api{OllamaApiClient}. Also implements \api{IChatClient}. \\
\pkg{Azure.AI.OpenAI} & Azure OpenAI client. \\
\pkg{Azure.Identity} & \api{DefaultAzureCredential} and friends. \\
\pkg{Microsoft.Extensions.Configuration.UserSecrets} & Keeping endpoints out of source control. \\
\pkg{Npgsql.EntityFrameworkCore.PostgreSQL} & Chapter~\ref{ch:memory}. \\
\pkg{OpenTelemetry.Exporter.OpenTelemetryProtocol} & Chapter~\ref{ch:observability}. \\
\bottomrule
\end{tabularx}
\end{center}

% ============================================================
\section{Keeping this appendix current}

The list above will be wrong within a quarter. Rather than trusting it, generate
your own:

\begin{shellcmd}
# Every package published by the Agent Framework account, newest first
curl -s "https://azuresearch-usnc.nuget.org/query?q=owner:MicrosoftAgentFramework&take=100&prerelease=true" \
  | jq -r '.data[] | "\(.id)\t\(.version)"' | sort
\end{shellcmd}

Run it before each chapter you work through, and diff against
\code{Directory.Packages.props}. Two minutes of scripting removes an entire class
of wasted afternoon.

\begin{note}
If you intend to publish anything derived from this book, automate the check in
CI: a weekly job that restores against the latest available versions and reports
build failures tells you which chapters have rotted before a reader does.
\end{note}
